\section{Linear transport and its distinguished coordinates}\label{sec:linear}
The Taylor estimates of Section~\ref{sec:norms} reduce the background flow to
its first four cumulant arrays.  We now identify the directions which linear
blocking does not contract.  There is one expanding mean direction and one
neutral covariance direction.  A spin insertion has one neutral direction
after multiplication by the scale factor $L^d$.  The distinction between a
covariance per unit area and a homogeneous perturbation is essential here.

\subsection{Coordinates and exact representatives}
Write $a(i)$ for the anchor cell of a slot $i\in \mathcal I_r$.  The array spaces and
indexed projective norms are those of Section~\ref{sec:blocking}: the tensor
norm at fixed indices is the projective tensor norm of the function sup norms,
and the array weight is $e^{v\sum n_i}$.  All unmarked arrays in this section
are invariant under grid translations and square symmetries.  A connected
array is zero unless the safety adjacency graph of its dependence carriers
is connected.  This is the support condition satisfied by joint cumulants of
the random lists.  Whenever indices are summed, an unused part of the
exponential weight pays the animal counts and polynomial size factors.

Let $\ell$ and $\ell_s$ be the charges of Lemma~\ref{ref:charges}.
For an even function and a scaled odd function, respectively, put
\[
 q_r(F)=r^{-1}\ell(F),\qquad q_{r,s}(U)=r^{-d}\ell_s(U).
\]
These are bounded functionals on a slot of length $n$, with norms at most
$C(n+S)$ and $C(n+S)^d$.  For a homogeneous mean array $m$ and a homogeneous
connected two-tensor $K$, define
\begin{align}
 x_r(m)&=\sum_{a(i)=c}q_r(m_i),\label{lin:x}\\
 g_r(K)&=\sum_{a(i)=c}\sum_{k\in \mathcal I_r}
                  (q_r\otimes q_r)(K_{i,k}).\label{lin:g}
\end{align}
The values do not depend on the cell $c$.  The second sum converges absolutely:
connectivity puts the second anchor within $C(n_i+n_k+S)r$ of the first,
and exponential weights dominate this number of placements.  The factor
$r^{-2}$ in~\eqref{lin:g} divides the charge covariance by the number of
microscopic sites in one cell.  Thus $g_r$ is a variance density.

Let $P_{Hr,c}$ denote reference conditional averaging to the square contour
with radius $Hr$ about the center of $c$, followed, for even functions, by
plane centering.  Define arrays supported on principal slots by
\begin{align}
 (v_j)_{p_c}&=r_j^{-1}\sum_{z\in c}P_{Hr_j,c}e(z),
       &x_{r_j}(v_j)&=1,\label{lin:v}\\
 (V_j)_{p_c,p_c}&=\sum_{z\in c}
                       \bigl(P_{Hr_j,c}e(z)\bigr)^{\otimes2},
       &g_{r_j}(V_j)&=1.\label{lin:V}
\end{align}
All other entries of $V_j$ vanish.  The charge identities follow from
$\ell(e(z))=1$ and $|c|=r_j^2$.  The bounds
$\|P_{Hr_j,c}e(z)\|\le C/r_j$ show that both representatives have uniformly
bounded indexed norms, including the projective norm of $V_j$.

Write $T_j$ for linear background transport and $T_j^{(k)}$ for its induced
transport on $k$-tensors, including the prescribed reindexing and sums.
Nested contours and the tower property give the exact identities
\begin{equation}\label{lin:exact}
 T_jv_j=Lv_{j+1},\qquad T_j^{(2)}V_j=V_{j+1},\qquad
 x_{r_{j+1}}T_j=Lx_{r_j},\qquad
 g_{r_{j+1}}T_j^{(2)}=g_{r_j}.
\end{equation}
For example, the $L^2$ child cells partition a parent cell, so transporting
\eqref{lin:v} yields
\[
 r_j^{-1}\sum_{z\in c'}P_{Hr_{j+1},c'}e(z).
\]
For~\eqref{lin:V}, each microscopic diagonal tensor is transported to the
same parent contour in both factors.  No off-diagonal microscopic pair is
created.  These observations also explain why the two multipliers differ.
The coordinate identities follow by preserving charges and dividing by
$R$ or $R^2$, respectively.  Reindexing long slots preserves the same total
charge sums.  We use the bounded projections
\[
 \Pi_{m,j}m=x_{r_j}(m)v_j,\qquad
 \Pi_{K,j}K=g_{r_j}(K)V_j.
\]

\begin{proposition}[Background linear splitting]\label{lin:background}
Fix the finite set of weights needed in Section~\ref{sec:norms}.  By taking
$L$ sufficiently large, all the following contraction constants can be made
smaller than any prescribed $\kappa>0$:
\[
 \|T_jm\|\le\kappa\|m\|\quad(x_{r_j}(m)=0),\qquad
 \|T_j^{(2)}K\|\le\kappa\|K\|\quad(g_{r_j}(K)=0),
\]
and $\|T_j^{(k)}A\|\le\kappa\|A\|$ for homogeneous connected arrays of
orders $k=3,4$.  Norms here and below have the fixed input and output weights
allocated to the corresponding estimate.  The bounds are uniform in $j$.
The two projections above intertwine transport with multipliers $L$ and $1$.
\end{proposition}

\begin{proof}
Put $M=(\log L)^3$.  If an input length exceeds $M$, the unused exponential
weight gives $C L^C e^{-cM}$ after all linear placement sums.  Charge
projection adds only polynomial factors in the lengths.  Consequently these
long inputs, including the charge needed to compensate their removal, have
norm at most $\eta_L\|A\|$, where $\eta_L\to0$ faster than every fixed
negative power of $L$.  It suffices to estimate the inputs of lengths at
most $M$; constants denoted by $C M^C$ include their weighted index sums.

For the mean, collect all short slots with anchor $c$ into one even
function supported in a square of radius $C M r$.  Homogeneity and square
symmetry make this grouped function square symmetric about $c$.  Its charge
is zero, up to the discarded long-input charge.  Subtract that charge using
a translated copy of $v_j$ before applying the reference estimate.  The
square-symmetric charge-zero exponent is $3-\nu$, so one child contribution
at the parent contour costs $C M^C L^{-3+\nu}$.  There are $L^2$ child
anchors.  Thus the short part is bounded by
$C M^C L^{-1+\nu}\|m\|$, plus $\eta_L\|m\|$ for the compensation.
The grouped function is formed before taking its sup norm; the triangle
inequality and the weighted slot sum bound that norm by the input array
norm.  This is where the square symmetry of the law is used.

For the covariance, connectivity and the short-length restriction put the
two anchors within $C M r$ of one another.  Separate pairs whose anchors go
to different parent cells from pairs going to the same parent.  In the first
case, an anchor lies in a strip of width $C M r$ about the boundary of a
parent cell.  There are $C L M$ such fine cells per parent.  Each centered
factor costs at most $C M/L$ under averaging; after the weighted sums over
slot shapes and the other anchor, this contribution is at most
\begin{equation}\label{lin:strip}
 C M^C L^{-1}\|K\|.
\end{equation}
This bound applies also to pairs inside one parent for which the first anchor
lies in that strip.

For the remaining pairs work first with one elementary projective tensor
$F\otimes G$.  If $c$ is its first anchor and $z_c$ the center site of $c$,
replace both factors by their charges times the same unit-charge field:
\[
 F\longmapsto q_r(F)\,rP_{Hr,c}e(z_c),\qquad
 G\longmapsto q_r(G)\,rP_{Hr,c}e(z_c).
\]
Their supports and the original supports are contained in a square of
radius $C M r$.  After subtracting the charge, the even improved estimate
and the tower property bound a replaced-factor error at the parent contour
by $C M^C L^{-2+\nu}$ times the original factor norm.  The other factor has
its usual $C M^C/L$ bound.  Summing $L^2$ first anchors gives
$C M^C L^{-1+\nu}\|K\|$ for terms containing a replacement error.

It remains to sum the term in which both factors were replaced.  With the
first anchor fixed, the transported elementary pair is the same for every
choice of its second slot or anchor.  Its scalar coefficient is therefore
the corresponding sum in~\eqref{lin:g}, and vanishes on the kernel of $g_r$.
For an anchor outside the strip, every short connected second anchor goes
to the same parent, so this cancellation is complete.  Restriction to short
inputs leaves only the discarded long-input charge, already bounded by
$\eta_L$.  Strip anchors were estimated in~\eqref{lin:strip}.  Homogeneity
is needed here to make the zero coordinate condition hold at every first
anchor; it is not a statement about an arbitrary nonhomogeneous covariance.
Taking infima over projective decompositions proves the asserted tensor
bound without any factor depending on the dimension of the function spaces.

Finally, a connected short $k$-tuple has all its anchors within $C_k M r$
of its first anchor.  The first anchor has $C L^2$ possible positions in a
rooted parent cell, and each of the $k$ centered factors costs $C M/L$.
Weighted sums bound the result by $C M^C L^{2-k}\|A\|$ for $k=3,4$.
Together with the long part and the preceding estimates this tends to zero.
The intertwining assertions follow from~\eqref{lin:exact}.
\end{proof}

\subsection{One marked point and tensor copies}
For a fixed marked site $y$, let $c_j(y)$ be its cell and put
\begin{equation}\label{lin:w}
 (w_{j,y})_{p_{c_j(y)}}=r_j^dP_{Hr_j,c_j(y)}\sigma_y.
\end{equation}
This is an odd scaled source list, with all other entries zero.  Its total
normalized spin charge is one.  For an arbitrary marked list $u$, and for
an array of $q$ marked copies whose slot indices are kept separate, set
\[
 Q_{j,y}(u)=\sum_i q_{r_j,s}(u_i),\qquad
 Q_{j,q,y}(U)=\sum_{i_1,\ldots,i_q}
                   q_{r_j,s}^{\otimes q}(U_{i_1,\ldots,i_q}).
\]
Every marked carrier contains $y$, so these are bounded functionals in the
weighted indexed spaces.  Write $\mathcal T_j=L^dT_j$ for scaled odd
transport and $\mathcal T_j^{(q)}$ for its tensor extension.  The exact
identities are
\[
 \mathcal T_jw_{j,y}=w_{j+1,y},\qquad
 Q_{j+1,q,y}\mathcal T_j^{(q)}=Q_{j,q,y}.
\]
The corresponding projection is
$\Pi_{q,j,y}U=Q_{j,q,y}(U)w_{j,y}^{\otimes q}$.

\begin{proposition}[Marked linear splitting]\label{lin:marked}
For each fixed finite set of positive integers $q$, the operators
$\mathcal T_j^{(q)}$ are bounded uniformly in $j,y,L$.  On
$\ker Q_{j,q,y}$ their norms tend to zero as $L\to\infty$.
A connected array comprising this marked tensor as one grouped argument and
at least one even background argument contracts without removal of a line.
The one-source kernel estimate also holds pathwise, and hence in each of the
weighted random-list $L^p$ norms used in Section~\ref{sec:norms}.
\end{proposition}

\begin{proof}
For a short source slot the ordinary odd estimate gives
$C((n+S)/L)^d$ before source scaling.  Multiplication by $L^d$ and summation
with the unused exponential weight give a constant independent of $L$;
one sums $(n+S)^d$ against that weight, rather than replacing it by $M^d$.
Long slots are again bounded by $C L^C e^{-cM}$.

For each short slot individually, subtract its normalized charge times the
representative at its input scale.  Both functions fit in a square of radius
$C(n+S)r$ containing $y$.  The charge-zero odd estimate, whose exponent is
$1+d-\nu$, followed by source scaling, bounds the difference between the
transported slot and its charged representative by
$C(n+S)^C L^{-1+\nu}$ times its norm.  Summing over slots is legitimate by
the exponential weights.  Long slots admit the same subtraction with a
beyond-polynomially small bound, since charge extraction has only a
polynomial size cost.  This gives a componentwise decomposition into a
bounded rank-one map and a remainder of arbitrarily small norm.

Apply that decomposition separately to each factor of an elementary
projective $q$-tensor.  The all-line term is exactly
$Q_{j,q,y}(U)w_{j+1,y}^{\otimes q}$.  Every other term contains a small
remainder and at most $q-1$ bounded factors.  The sum of its $2^q-1$ choices
is small for fixed $q$.  Infimizing over decompositions proves the result
in the projective norm.  In particular no norm comparison involving the
number of spin configurations is used.

A connected array with an additional background factor has no $L^2$ volume
sum: its root is the fixed mark.  Connectivity bounds the remaining anchor
sums by powers of the slot lengths.  At least one even factor supplies
$C M^C/L$, whereas the scaled marked factors stay bounded.  This proves
contraction in the mixed sector.  All these are componentwise deterministic
inequalities, so applying the center-coordinate projection to each random
realization and then taking its $L^p$ norm gives the last assertion.  The
buffer $H$ keeps every point of a cell uniformly away from its cutting
contour, making the constants independent of the position of $y$ in its
cell.
\end{proof}
